\scanpage{20}
% lines 1
\[
f'(X)=2\widetilde X,\ f'(Y)=2(\widetilde Y),\ f'(H)=\widetilde H \tag*{(66)}
\]

% lines 2
La forme bilinéaire \deleted{associée} \inserted{$\varphi'$ correspondante sur $E'$} est donc définie par

% lines 3–4
\[
\left\{\begin{aligned}
\varphi'(\deleted{\unsure{\widetilde X}}X,X)=\varphi'(Y,Y)&=0 &\varphi'(H,H)&=1\\
\varphi'(X,\deleted{\unsure{Y}}H)=\varphi'(Y,H)&=0 &\varphi'(X,Y)&=2
\end{aligned}\right. \tag*{(67)}
\]

% lines 5–7
\[
\begin{aligned}
&\text{i.e. en termes de \unsure{v|$\eta$}}\\
\eta&=x\overset{\unsure{\leftrightarrow|\frown}}{X+h}H+yY\\
\eta'&=x'X+h'H+y'Y
\end{aligned}
\]

% lines 8
la forme $\varphi'$ s'explicite par

% lines 9
\[
\varphi'(\eta,\eta')=\varphi((x,h,y);(x',h',y'))=hh'+2(xy'+yx') \tag*{(68)}
\]

% lines 10–11
la forme quadratique associée $q'$ étant donnée par

% lines 12
\[
q'(\eta)=\varphi'(\eta,\eta)=h^2+4xy \tag*{(69)}
\]

% lines 13–16
\deleted{``En \unsure{général}''} \inserted{Si 2 non inv dans $S$} $\varphi'$ n'est pas associée à une forme quadratique, \deleted{en particulier en car.} car en car. 2 la forme $q'(\eta)$ se réduit à $h^2$, qui est non nulle\ldots

% lines 16–25
On notera que dans les coordonnées universelles choisie, permettant d'exprimer $q,\varphi,q',\varphi'$ \deleted{en} termes de \inserted{quatre} expressions \unsure{polynômiales|polynomiales} à coefficients entiers, aucune de ces expressions n'est divisible par un entier $>1$ -- i.e. les coeff. sont premiers entre eux. Cela signifi aussi que \deleted{au dessus de \unsure{n'importe} \unsure{ou}} \inserted{si $S$ est le spectre d'un corps} (de car. quelc.), ces quatre fonctions sont $\ne0$.

% lines foot
% archive sheet number: 19
